\section{发布与评测手册：从离线验证到线上灰度}

讲者反复强调 Agent 是 non-deterministic software，因此发布流程不能照搬传统 Web 接口。这里给出一个与课程观点一致的发布手册，可直接映射到团队周度节奏。

\subsection{四层发布门禁（Release Gates）}

\begin{table}[H]
\centering
\caption{建议的 Agent 发布门禁}
\begin{tabularx}{\textwidth}{>{\raggedright\arraybackslash}p{2.2cm} >{\raggedright\arraybackslash}p{3.6cm} >{\raggedright\arraybackslash}X}
\toprule
门禁层级 & 核心检查项 & 通过标准示例 \\
\midrule
Gate 1: 静态检查 & 策略规则完整性、工具 schema 兼容性、敏感词与动作白名单 & 无 blocking 规则冲突，schema 校验 100\% 通过 \\
Gate 2: 仿真回归 & Towbench 类多轮任务、工具调用一致性、策略遵循率 & 关键场景成功率达到基线，策略违规率低于阈值 \\
Gate 3: 影子流量 & 与人工或旧版本并行运行，不对用户生效 & 行为差异可解释，重大风险事件为 0 \\
Gate 4: 小流量灰度 & 真实用户 1\%-5\% 流量，实时监控升级率与失败类型 & 升级率不劣于基线，安全告警可控，可随时一键回滚 \\
\bottomrule
\end{tabularx}
\end{table}

\begin{importantbox}{发布纪律的底线}
任何会改写系统 of record 的 Agent 版本，都必须满足“可回滚、可审计、可解释”三要素。没有回滚开关的发布，在客服域等同于无保护上线。
\end{importantbox}

\subsection{评测集设计：覆盖而非堆量}

许多团队把评测等同于“收集更多数据”。更有效做法是按风险分层设计评测集：
\begin{itemize}[leftmargin=1.5em]
    \item 高频低风险：查询、进度解释、标准政策说明；
    \item 高频高风险：退款、改址、账户权限变更；
    \item 低频高风险：合规边界问题、异常账户、争议升级；
    \item 对抗样本：prompt injection、多语种混合、错别字、口音扰动。
\end{itemize}

\begin{knowledgebox}{“覆盖率”建议定义}
一个可操作的覆盖率定义可以是：
\[
\text{Coverage} = \frac{\text{被评测集触达的关键策略分支数}}{\text{生产环境关键策略分支总数}}
\]
比起单纯增加样本量，更应关注关键策略分支是否被触达。
\end{knowledgebox}

\subsection{线上监控与告警策略}

灰度期必须同时监控业务指标、质量指标、安全指标。建议最小看板包含：
\begin{itemize}[leftmargin=1.5em]
    \item 业务：会话量、可验证解决率、单位会话成本；
    \item 质量：P50/P90 响应延迟、升级率、重复联系率；
    \item 安全：输入攻击命中率、输出越界命中率、人工紧急接管率。
\end{itemize}

\begin{warningbox}{只看均值是高频事故源}
如果只看平均延迟和平均成功率，容易错过高分位尾部风险。客服系统中的投诉和舆情往往由少量极差会话触发，因此必须持续追踪 P95/P99 与长尾失败模式。
\end{warningbox}

\subsection{本章小结}

Agent 发布应从“功能上线”升级为“风险受控上线”。分层门禁、风险分层评测和高分位监控是三件必须同时具备的基本功。

